% GENERATED FILE. Edit canonical lessons, then run npm run generate:books.
% canonical-source-hash: fnv1a64:f4241b6a9b28ce24
% canonical-lessons: JA-C84-donna, JA-C84-hoka, JA-C84-minna, JA-C84-kouiu, JA-C84-aaiu

\chapter{What kind of, Other, Everybody, This sort of, That sort of}
\label{ch:ja-what-kind-of-other-everybody-this-sort-of-that-sort-of}
\hlchaptermodality{84}

\begin{chapteropening}
\textbf{By the end of this chapter:} \emph{I can say what kind of, other, everybody, this sort of and that sort of.}

Complete the last lesson of chapter 84: i can say what kind of, other, everybody, this sort of and that sort of.
\end{chapteropening}

\section[\texorpdfstring{\ja{どんな} (donna)}{どんな (donna)}]{\ja{どんな} (donna) — what kind of}

\label{lesson:JA-C84-donna}



\begin{quote}
Before the new one: say the Japanese for this kind of, then the Japanese for that kind of.
\end{quote}

\subsection*{You'll want to know: \ja{どんな}}
\textbf{\ja{どんな}} — \emph{donna} — \textquotedblleft{}what kind of\textquotedblright{}.

\begin{grammarlens}[title={What you've built}]
One more word for this chapter.
\end{grammarlens}

\subsection*{Your turn}
\begin{itemize}
\raggedright
  \item \emph{Say it:} \emph{donna}
  \item \emph{Say it:} \emph{donna}, once more
  \item \emph{Say it:} say \emph{donna}
  \item \emph{Recall:} say the Japanese for this kind of, then the Japanese for that kind of, then say \emph{donna} again
\end{itemize}

\begin{tcolorbox}[breakable,colback=teal!4,colframe=teal!35!black,title={Before you move on}]
What does \ja{どんな} mean? (\textquotedblleft{}What kind of\textquotedblright{}.) Say it once more.
\end{tcolorbox}
\section[\texorpdfstring{\ja{ほか} (hoka)}{ほか (hoka)}]{\ja{ほか} (hoka) — other, another}

\label{lesson:JA-C84-hoka}



\begin{quote}
Before the new one: say the Japanese for that kind of, then the Japanese for what kind of.
\end{quote}

\subsection*{You'll want to know: \ja{ほか}}
\textbf{\ja{ほか}} — \emph{hoka} — \textquotedblleft{}other, another\textquotedblright{}.

\begin{grammarlens}[title={What you've built}]
One more word for this chapter.
\end{grammarlens}

\subsection*{Your turn}
\begin{itemize}
\raggedright
  \item \emph{Say it:} \emph{hoka}
  \item \emph{Say it:} \emph{hoka}, once more
  \item \emph{Say it:} say \emph{hoka}
  \item \emph{Recall:} say the Japanese for that kind of, then the Japanese for what kind of, then say \emph{hoka} again
\end{itemize}

\begin{tcolorbox}[breakable,colback=teal!4,colframe=teal!35!black,title={Before you move on}]
What does \ja{ほか} mean? (\textquotedblleft{}Other, another\textquotedblright{}.) Say it once more.
\end{tcolorbox}
\section[\texorpdfstring{\ja{みんな} (minna)}{みんな (minna)}]{\ja{みんな} (minna) — everybody}

\label{lesson:JA-C84-minna}



\begin{quote}
Before the new one: say the Japanese for what kind of, then the Japanese for other.
\end{quote}

\subsection*{You'll want to know: \ja{みんな}}
\textbf{\ja{みんな}} — \emph{minna} — \textquotedblleft{}everybody\textquotedblright{}.

\begin{grammarlens}[title={What you've built}]
One more word for this chapter.
\end{grammarlens}

\subsection*{Your turn}
\begin{itemize}
\raggedright
  \item \emph{Say it:} \emph{minna}
  \item \emph{Say it:} \emph{minna}, once more
  \item \emph{Say it:} say \emph{minna}
  \item \emph{Recall:} say the Japanese for what kind of, then the Japanese for other, then say \emph{minna} again
\end{itemize}

\begin{tcolorbox}[breakable,colback=teal!4,colframe=teal!35!black,title={Before you move on}]
What does \ja{みんな} mean? (\textquotedblleft{}Everybody\textquotedblright{}.) Say it once more.
\end{tcolorbox}
\section[\texorpdfstring{\ja{こういう} (kouiu)}{こういう (kouiu)}]{\ja{こういう} (kouiu) — this sort of}

\label{lesson:JA-C84-kouiu}



\begin{quote}
Before the new one: say the Japanese for other, then the Japanese for everybody.
\end{quote}

\subsection*{You'll want to know: \ja{こういう}}
\textbf{\ja{こういう}} — \emph{kouiu} — \textquotedblleft{}this sort of\textquotedblright{}.

\begin{grammarlens}[title={What you've built}]
One more word for this chapter.
\end{grammarlens}

\subsection*{Your turn}
\begin{itemize}
\raggedright
  \item \emph{Say it:} \emph{kouiu}
  \item \emph{Say it:} \emph{kouiu}, once more
  \item \emph{Say it:} say \emph{kouiu}
  \item \emph{Recall:} say the Japanese for other, then the Japanese for everybody, then say \emph{kouiu} again
\end{itemize}

\begin{tcolorbox}[breakable,colback=teal!4,colframe=teal!35!black,title={Before you move on}]
What does \ja{こういう} mean? (\textquotedblleft{}This sort of\textquotedblright{}.) Say it once more.
\end{tcolorbox}
\section[\texorpdfstring{\ja{ああいう} (aaiu)}{ああいう (aaiu)}]{\ja{ああいう} (aaiu) — that sort of}

\label{lesson:JA-C84-aaiu}



\begin{quote}
Before the new one: say the Japanese for everybody, then the Japanese for this sort of.
\end{quote}

\subsection*{You'll want to know: \ja{ああいう}}
\textbf{\ja{ああいう}} — \emph{aaiu} — \textquotedblleft{}that sort of\textquotedblright{}.

\begin{grammarlens}[title={What you've built}]
One more word for this chapter.
\end{grammarlens}

\subsection*{Your turn}
\begin{itemize}
\raggedright
  \item \emph{Say it:} \emph{aaiu}
  \item \emph{Say it:} \emph{aaiu}, once more
  \item \emph{Say it:} say \emph{aaiu}
  \item \emph{Recall:} say the Japanese for everybody, then the Japanese for this sort of, then say \emph{aaiu} again
\end{itemize}

\begin{tcolorbox}[breakable,colback=teal!4,colframe=teal!35!black,title={Before you move on}]
What does \ja{ああいう} mean? (\textquotedblleft{}That sort of\textquotedblright{}.) Say it once more.
\end{tcolorbox}
